%% ma: L12-B14-0133
%% lop: 12
%% chuong: 5
%% bai: 14
%% dang: 12.14.30
%% loai: TN4
%% muc: VD
%% dap_an: "D"
%% nguon: "Ảnh giáo viên gửi 25/11/2026 (tài liệu ôn tập Phương trình mặt phẳng), Phần I Câu 24"
%% kien_thuc: "Bài 14 (phương trình mặt phẳng); tứ diện OABC có OA,OB,OC đôi một vuông góc thì OH⊥(ABC) với H là trực tâm tam giác ABC"
%% trang_thai: chinh-thuc
\begin{ex}
Trong không gian $Oxyz$, cho điểm $H(1;2;-2)$. Mặt phẳng $(\alpha)$ đi qua $H$ và cắt các trục $Ox,Oy,Oz$ lần lượt tại các điểm $A,B,C$ sao cho $H$ là trực tâm của tam giác $ABC$ có phương trình là
\choice{$y+z=0$}{$2x+y-2z-8=0$}{$x+2y+2z-1=0$}{\True $x+2y-2z-9=0$}
\loigiai{Vì $OA,OB,OC$ đôi một vuông góc và $H$ là trực tâm tam giác $ABC$ nên $OH\perp(ABC)$. Do đó $\overrightarrow{OH}=(1;2;-2)$ là một vectơ pháp tuyến của $(\alpha)$.
Phương trình $(\alpha)$: $(x-1)+2(y-2)-2(z+2)=0\Leftrightarrow x+2y-2z-9=0$.}
\end{ex}
